%%%PAGE 142 rot=0 legibility=good summary=卫星几何类习题：同步卫星微波信号传播时间(6)、赤道上连续看到卫星的时间(7)、A/B卫星不能直接通讯的时间(8)
%%%BLOCK id=142-01 ch=05 sec=同步卫星·信号传播 kind=problem alt=none cont=none y=0.03-0.46
\begin{problem}[6]
\textbf{6、}我国同步卫星、定点位置与东经98$^\circ$经线在同一平面内。把甘肃省嘉峪关处的经度和纬度近似为东经98$^\circ$和北纬$\alpha=40^\circ$，地球半径$R$，自转周期$T$，地表重力加速度$g$，光速$c$，求该同步卫星发出微波信号传到嘉峪关所需时间。% 原稿中“北纬α=40°”被圈出
%FIG p142_a 0.17 0.185 0.45 0.32
\notefig[0.3]{figures/p142_a.png}
\begin{solution}
\[ l=\sqrt{r^2+R^2-2rR\cos\alpha} \]
由$\dfrac{GMm}{r^2}=m\dfrac{4\pi^2}{T^2}r$　和$gR^2=GM$
\[ \therefore\ r=\sqrt[3]{\dfrac{gR^2T^2}{4\pi^2}} \]
\[ t=\dfrac{l}{c} \]
解出
\[ t=\dfrac{\sqrt{\left(\dfrac{gR^2T^2}{4\pi^2}\right)^{\frac{3}{2}}+R^2-2R\left(\dfrac{gR^2T^2}{4\pi^2}\right)^{\frac{1}{3}}\cos 40^\circ}}{c} \] % CHECK: 第一项指数手写为 3/2，按 r^2 应为 2/3（照抄）
\end{solution}
\end{problem}
%%%END
%%%BLOCK id=142-02 ch=05 sec=追及相遇·卫星过顶 kind=problem alt=none cont=none y=0.46-0.83
\begin{problem}[7]
\textbf{7、}地球自转周期$T_0$，在赤道上一点的人能连续看到该卫星的时间是多少？（在$B_1$的人刚看到$A$卫星）
%FIG p142_b 0.10 0.465 0.39 0.64
\notefig[0.3]{figures/p142_b.png}
\begin{solution}
\[ \omega t-\omega_0 t=\dfrac{2\pi}{3} \]
\[ \left(\dfrac{2\pi}{T}-\dfrac{2\pi}{T_0}\right)t=\dfrac{2\pi}{3} \]
\[ \begin{cases}
\text{\struck{$R+$}}\,r\,\text{\struck{\unclear{$\ast h$}}}=2R\\[4pt]
\dfrac{GMm}{r^2}=m\dfrac{4\pi^2}{T^2}r
\end{cases} \] % 第一式左侧有划去的字迹，划去后读作 r=2R
\[ T_{\text{卫}}=2\pi\sqrt{\dfrac{r^3}{GM}}=2\pi\sqrt{\dfrac{8R^3}{gR^2}}=4\pi\sqrt{\dfrac{2R}{g}} \] % CHECK: T 的下标手写似“工”，应为“卫”
\noindent\struck{$=2\pi\sqrt{\dfrac{r}{g}}=2\pi\sqrt{\dfrac{2R}{g}}$} % 原稿此行被划去
\[ \therefore\ t=\dfrac{4\pi\sqrt{\dfrac{2R}{g}}\,T}{3\left(T_0-4\pi\sqrt{\dfrac{2R}{g}}\right)} \] % CHECK: 分子末尾的 T 手写无下标，按推导应为 T_0
\end{solution}
\end{problem}
%%%END
%%%BLOCK id=142-03 ch=05 sec=追及相遇·卫星通讯 kind=problem alt=none cont=none y=0.83-1.00
\begin{problem}[8]
\textbf{8、}自转方向与地球自转方向一样\\
地球自转周期$T$和$R$，(1)求$B$的周期　(2)$A$、$B$不能直接通讯的时间% “直接”为行上插入（有脱字符）
%FIG p142_c 0.045 0.872 0.335 0.985
\notefig[0.3]{figures/p142_c.png}
\begin{solution}
(1)　由开三　$\dfrac{r^3}{T_B^2}=\dfrac{h^3}{T^2}$
\[ T_B=T\sqrt{\dfrac{r^3}{h^3}} \]
\unclear{sin}% 此处(1)与(2)之间有零星字迹“sin”
(2)　$\omega_B t-\omega t=2\theta$
\[ \sin\beta=\dfrac{R}{r},\qquad \sin\alpha=\dfrac{R}{h},\qquad \theta=\alpha+\beta \]
\[ \therefore\ t= \]
\end{solution}
\end{problem}
%%%END
%%%PAGEEND 142
